% barrier_groups.png: the four groups of jumps around a cluster, drawn on a triangular lattice. The sites and the jumps come from
% analysis_schematics.py. Make it in this folder with, for any scratch folder T:
%   python3 analysis_schematics.py T
%   TEXINPUTS=.:T: pdflatex -output-directory=T fig_barrier_groups.tex
%   python3 pdf_to_png.py T/fig_barrier_groups.pdf ../barrier_groups.png 1100
\documentclass[11pt,tikz,border=1.5pt]{standalone}
% Shared preamble of all figures. Every fig_*.tex is compiled on its own, see the comment at its top; the fonts are those of the matplotlib
% style xzc_sans (amsmath, amssymb, amsfonts, sfmath), so that line drawings and data plots look the same.
\usepackage{amsmath,amssymb,amsfonts}
\usepackage{sfmath}
\renewcommand{\familydefault}{\sfdefault}
\usepackage{graphicx}
\usetikzlibrary{arrows.meta,calc,positioning,decorations.pathreplacing,backgrounds,shapes.geometric}
% species: Al gray, Mg orange, Zn blue; red marks the vacancy and the jumping atom
\definecolor{cAl}{HTML}{B8BCC2}
\definecolor{cMg}{HTML}{E08A2E}
\definecolor{cZn}{HTML}{3F78B5}
\definecolor{cOne}{HTML}{C8423B}
\definecolor{cTwo}{HTML}{3F78B5}
\definecolor{cThree}{HTML}{3E9B5C}
\definecolor{cFour}{HTML}{E0A21B}
\definecolor{cGray}{HTML}{7E858F}
\definecolor{cDark}{HTML}{2B2B2B}
\newcommand{\panel}[1]{\textbf{(#1)}}
\newcommand{\DM}{D_{\mathrm{MEP}}}
\newcommand{\KI}{K_{\mathrm{I}}}
\newcommand{\KF}{K_{\mathrm{F}}}
\newcommand{\dE}{\Delta E}
\newcommand{\Ea}{\Delta E_{\mathrm{a}}}
\newcommand{\ea}{\Delta e_{\mathrm{a}}}
\newcommand{\Ebind}{E_{\mathrm{bind}}^{\mathrm{site}}}
% a ball with a soft highlight, like the rendered spheres
\makeatletter
\pgfdeclareradialshading[tikz@ball]{softball}{\pgfpoint{-9bp}{9bp}}{%
  color(0bp)=(tikz@ball!22!white); color(6bp)=(tikz@ball!62!white); color(13bp)=(tikz@ball!96!white);
  color(20bp)=(tikz@ball!92!black); color(27bp)=(tikz@ball!74!black); color(50bp)=(tikz@ball!55!black)}
\makeatother
\tikzset{
  every picture/.append style={line cap=round, line join=round},
  every node/.append style={font=\small},
  img/.style={inner sep=0pt, anchor=south west},
  lab/.style={font=\normalsize, inner sep=1pt, anchor=north west},
  above lab/.style={font=\normalsize, inner sep=1pt, anchor=south west},
  capt/.style={align=center, anchor=north},
  atom/.style={circle, ball color=#1, shading=softball, minimum size=12pt, inner sep=0pt},
  atom/.default=cAl,
  vac/.style={circle, draw=cOne, line width=1.1pt, fill=white, minimum size=11.5pt, inner sep=0pt},
  arr/.style={-{Stealth[length=5pt]}, line width=0.8pt},
  dbl/.style={{Stealth[length=4pt]}-{Stealth[length=4pt]}, line width=0.8pt},
  jump/.style={arr, cOne, line width=1.2pt},
  block/.style={-{Triangle[length=8pt, width=13pt]}, line width=5pt, cTwo!38},
  cell axis/.style={-{Triangle[length=6pt, width=5.5pt]}, line width=1.2pt, black},
  dir/.style={fill=white, fill opacity=0.75, text opacity=1, font=\footnotesize, inner sep=1.5pt},
}
\graphicspath{{img/}}
% data plots are included at one scale: \plot{name}{x}{y}{file}
\newcommand{\plotscale}{0.6}
\newcommand{\plot}[4]{\node[img] (#1) at (#2,#3) {\includegraphics[scale=\plotscale]{#4}};}
\newcommand{\plotlabel}[2]{\node[above lab] at ($(#1.north west)+(-0.05,0.02)$) {\panel{#2}};}
% legend of the species, first entry at (#1,#2): Al, Mg, Zn (and the vacancy)
\tikzset{leg/.style={anchor=base west, inner sep=2pt, yshift=-3.2pt}}
\newcommand{\specieslegend}[2]{%
  \foreach \lx/\lc/\lt in {0/cAl/Al, 1.2/cMg/Mg, 2.4/cZn/Zn} {
    \node[atom=\lc, minimum size=10pt] at ({#1+\lx},#2) {}; \node[leg] at ({#1+\lx+0.2},#2) {\lt};}}
\newcommand{\vacancylegend}[2]{\node[vac, minimum size=9.5pt] at (#1,#2) {}; \node[leg] at ({#1+0.2},#2) {vacancy};}

\input{barrier_groups_sites}
\begin{document}
\begin{tikzpicture}[x=0.92cm, y=0.92cm,
    al/.style={atom=cAl, minimum size=13.4pt}, mg/.style={atom=cMg, minimum size=15.4pt}, zn/.style={atom=cZn, minimum size=13.4pt},
    rim/.style={cTwo!75, line width=27.5pt}, region/.style={cTwo!20, line width=25.9pt},
    near rim/.style={cGray!75, line width=27.5pt}, near region/.style={cGray!18, line width=25.9pt},
    one/.style={-{Stealth[length=8pt, width=7.6pt]}, line width=2.3pt, shorten <=3pt, shorten >=1pt},
    both/.style={{Stealth[length=8pt, width=7.6pt]}-{Stealth[length=8pt, width=7.6pt]}, line width=2.3pt, shorten <=1pt, shorten >=1pt},
    swatch/.style={line width=0.8pt, rounded corners=3pt}]
  \newcommand{\lines}[2]{\foreach \x/\y/\u/\v in #1 \draw[#2] (\x,\y) -- (\u,\v);}
  \newcommand{\nodes}[2]{\foreach \x/\y in #1 \node[#2] at (\x,\y) {};}
  \newcommand{\jumps}[3]{\lines{#1}{#2, white, line width=4.2pt} \lines{#1}{#2, #3}}
  \lines{\bondsNear}{near rim} \lines{\bondsNear}{near region}
  \lines{\bondsCluster}{rim} \lines{\bondsCluster}{region}
  \nodes{\sitesAl}{al} \nodes{\clusterMg}{mg} \nodes{\clusterZn}{zn}
  \jumps{\jumpsIn}{both}{cFour}
  \jumps{\jumpsTo}{one}{cOne}
  \jumps{\jumpsOn}{both}{cThree}
  \jumps{\jumpsOff}{one}{cTwo}
  % ---------------- legend ----------------
  \begin{scope}[shift={(7.4,0.55)}, x=1cm, y=1cm]
    \foreach \t/\s/\c [count=\i from 0] in {in/both/cFour, to/one/cOne, on/both/cThree, off/one/cTwo} {
      \draw[\s, \c] (0,-0.7*\i) -- ++(0.92,0); \node[leg] at (1.15,-0.7*\i) {\texttt{\t}};}
    \foreach \t/\c/\f [count=\i from 0] in {cluster/cTwo/20, nearest neighbors/cGray/18} {
      \draw[swatch, \c!75, fill=\c!\f] (0.1,-3.30-0.75*\i) rectangle ++(0.72,0.46); \node[leg] at (1.15,-3.07-0.75*\i) {\t};}
    \foreach \t [count=\i from 0] in {Al, Mg, Zn} {
      \node[atom=c\t, minimum size=12pt] at (0.22+1.18*\i,-4.75) {}; \node[leg] at (0.46+1.18*\i,-4.75) {\t};}
  \end{scope}
\end{tikzpicture}
\end{document}
